\documentclass[11pt,a4paper]{article}
\usepackage{fontspec,amsmath,amssymb,amsthm,microtype}
\setmainfont{lmroman10-regular.otf}[BoldFont=lmroman10-bold.otf,ItalicFont=lmroman10-italic.otf,BoldItalicFont=lmroman10-bolditalic.otf]
\usepackage{xcolor}
\usepackage[most]{tcolorbox}
\usepackage{mathrsfs}
\providecommand{\lt}{<}
\providecommand{\gt}{>}
\setmonofont{DejaVuSansMono.ttf}[BoldFont=DejaVuSansMono-Bold.ttf,ItalicFont=DejaVuSansMono-Oblique.ttf,BoldItalicFont=DejaVuSansMono-BoldOblique.ttf,Scale=0.86]
\newfontfamily\leanfallback{STIXTwoMath-Regular.otf}[Scale=0.86]
\definecolor{leanrule}{gray}{0.62}
\definecolor{leanback}{gray}{0.962}
\definecolor{leantitle}{gray}{0.22}
\definecolor{leanmuted}{gray}{0.38}
\newtcolorbox{leanbox}[1]{enhanced,breakable,colback=leanback,colframe=leanrule,
  boxrule=0.4pt,arc=0.6mm,left=2.4mm,right=2.4mm,top=1.8mm,bottom=1.8mm,
  before skip=0.8em,after skip=1em,fonttitle=\small\bfseries,coltitle=leantitle,
  colbacktitle=leanback,titlerule=0.3pt,title={#1}}
\newlength{\leancharwidth}
\newenvironment{leancode}{\par\smallskip\begingroup\ttfamily\footnotesize\raggedright
  \settowidth{\leancharwidth}{M}\setlength{\parskip}{0pt}\setlength{\parindent}{0pt}
  \frenchspacing}{\par\endgroup\smallskip}
\newcommand{\leanline}[2]{\hangindent=\dimexpr#1\leancharwidth+4\leancharwidth\relax
  \hangafter=1\noindent\hspace*{#1\leancharwidth}#2\par}
\newcommand{\leansource}[1]{\par{\footnotesize\color{leanmuted}#1}\par}
\newcommand{\leanlabel}[1]{\par\smallskip\noindent{\small\bfseries #1}\par}
\usepackage[margin=25mm]{geometry}
\usepackage[unicode,hidelinks]{hyperref}
\hypersetup{pdfauthor={LangevinLean Project}}
\setlength{\emergencystretch}{3em}
\setlength{\parskip}{0.4em}
\setlength{\parindent}{0pt}
\allowdisplaybreaks[2]
\raggedbottom
\title{\texorpdfstring{Wasserstein Mixing Time of the Unadjusted Langevin Algorithm}{Wasserstein Mixing Time of the Unadjusted Langevin Algorithm}}
\author{LangevinLean Project}
\date{}
\begin{document}\maketitle
Draft mathematical exposition; these notes may contain mistakes.
The exact source bias, finite-time and accuracy bounds are proved under the original \(C^2\) Hessian assumptions.
\section{\texorpdfstring{Theorem}{Theorem}}


Let \(V\in C^2(\mathbb R^d)\), \(0<\alpha\le\beta\), and
\(\alpha I_d\preccurlyeq\nabla^2V\preccurlyeq\beta I_d\).
For \(0<h\le1/\beta\), set \(\kappa=\beta/\alpha\).
The Gibbs target is \(\pi(dx)=Z^{-1}e^{-V(x)}dx\), and
\(\widehat\pi\) is the constructed invariant law of the actual independent-Gaussian ULA transition. Then


\[W_2(\pi,\widehat\pi)\le6h\sqrt{\kappa\beta d}.\]


The invariant law, stationary Langevin process, reversible Markov semigroup,
Poisson corrector and synchronous comparison are all constructed from these
source assumptions. No spectral, Markov, invariant-law or comparison estimate
is an additional hypothesis of this endpoint.

\begin{leanbox}{Lean: \texttt{bias\_\allowbreak{}bound}}
\leansource{\texttt{bias\_\allowbreak{}bound} \quad WassersteinMixing/\allowbreak{}Theorem.lean, lines 21–27}
\begin{leancode}
\leanline{0}{theorem\ bias\_bound\ (P\ :\ Parameters)\ (V\ :\ E\ P.dimension\ →\ ℝ)}
\leanline{4}{(hA\ :\ HessianAssumptions\ P\ V)\ :}
\leanline{4}{ULA.wasserstein2Real\ (ULA.gibbsMeasure\ V)}
\leanline{6}{(invariantLaw\ P\ V\ hA.toPotentialAssumptions)\ ≤}
\leanline{8}{6\ *\ P.step\ *\ Real.sqrt\ (P.condition\ *\ P.beta\ *\ P.dimension)}
\end{leancode}
\end{leanbox}

\section{\texorpdfstring{Deriving the potential contract from the source Hessian bounds}{Deriving the potential contract from the source Hessian bounds}}


The source assumes precisely \(V\in C^2\) and
    \(\alpha I_d\preccurlyeq\nabla^2V\preccurlyeq\beta I_d\).
    Symmetry and the quadratic-form bounds give
    \(\|\nabla^2V(x)\|\le\beta\). The mean-value theorem makes
    \(\nabla V\) globally \(\beta\)-Lipschitz. Integrating the lower Hessian
    inequality along line segments gives the strong supporting-plane
    inequality, hence \(\alpha\)-strong convexity.




Thus the first-order contract used by the existing Langevin library is
    a consequence of the paper's original assumptions. The final proof works
    directly under \(C^2\); it does not require the source proof's temporary
    smoothness and polynomial-growth assumption or a smoothing limit.

\begin{leanbox}{Lean: \texttt{HessianAssumptions.\allowbreak{}toPotentialAssumptions}}
\leansource{\texttt{HessianAssumptions.\allowbreak{}toPotentialAssumptions} \quad WassersteinMixing/\allowbreak{}CurvatureAssumptions.lean, lines 138–142}
\begin{leancode}
\leanline{0}{theorem\ HessianAssumptions.toPotentialAssumptions\ (P\ :\ Parameters)\ (V\ :\ H\ →\ ℝ)}
\leanline{4}{(hV\ :\ HessianAssumptions\ P\ V)\ :\ PotentialAssumptions\ P\ V}
\end{leancode}
\leanlabel{How the Lean proof follows the argument}
\texttt{Hessian\allowbreak{}Assumptions} records the quadratic forms of the
    genuine derivative of the gradient. Its
    \texttt{to\allowbreak{}Potential\allowbreak{}Assumptions} theorem derives every field of the
    internal potential contract.
\end{leanbox}

\section{\texorpdfstring{Constructing the actual stationary Brownian comparison}{Constructing the actual stationary Brownian comparison}}


Start the diffusion at an actual Gibbs-distributed random vector
    independent of a canonical Brownian path. Picard iteration in the
    complete \(L^2\) space constructs the global solution of


    \[X_t-X_0=-\int_0^t\nabla V(X_s)\,ds+\sqrt2\,B_t.\]


Uniqueness identifies its successive segments with the reference
    process used by the synchronous Gaussian ULA chain. The Gibbs law is
    stationary. The independent Gaussian increments have the exact laws
    required by the original ULA recursion; their pushforward chain laws are
    the actual iterates of that transition.

\begin{leanbox}{Lean: \texttt{state\_\allowbreak{}integral}, \texttt{state\_\allowbreak{}law}, \texttt{actualChain\_\allowbreak{}law}}
\leansource{\texttt{state\_\allowbreak{}integral} \quad WassersteinMixing/\allowbreak{}SemigroupGlobal.lean, lines 133–150}
\begin{leancode}
\leanline{0}{theorem\ state\_integral\ (P\ :\ Parameters)\ (V\ :\ E\ P.dimension\ →\ ℝ)}
\leanline{4}{(hV\ :\ PotentialAssumptions\ P\ V)\ (t\ :\ ℝ)\ (ht\ :\ 0\ ≤\ t)\ :}
\leanline{4}{letI\ :=\ actualProbability\_isProbability\ P\ V\ hV}
\leanline{4}{state\ P\ V\ hV\ t\ -\ initial\ P\ V\ hV\ =}
\leanline{6}{-(∫\ s\ in\ 0..t,\ ULA.lpDrift\ (actualProbability\ P\ V)\ (gradient\ V)}
\leanline{8}{hV.gradientLipschitz\ (state\ P\ V\ hV\ s))\ +}
\leanline{6}{ULA.brownianLangevinNoise\ (actualProbability\ P\ V)\ (brownian\ P)}
\leanline{8}{(brownian\_preBrownian\ P\ V\ hV)\ 0\ t}
\end{leancode}
\leansource{\texttt{state\_\allowbreak{}law} \quad WassersteinMixing/\allowbreak{}SemigroupGlobal.lean, lines 206–217}
\begin{leancode}
\leanline{0}{theorem\ state\_law\ (P\ :\ Parameters)\ (V\ :\ E\ P.dimension\ →\ ℝ)}
\leanline{4}{(hV\ :\ PotentialAssumptions\ P\ V)\ (t\ :\ ℝ)\ :}
\leanline{4}{(actualProbability\ P\ V).map\ (state\ P\ V\ hV\ t)\ =\ ULA.gibbsMeasure\ V}
\end{leancode}
\leansource{\texttt{actualChain\_\allowbreak{}law} \quad WassersteinMixing/\allowbreak{}ActualComparison.lean, lines 193–212}
\begin{leancode}
\leanline{0}{theorem\ actualChain\_law\ (P\ :\ Parameters)\ (V\ :\ E\ P.dimension\ →\ ℝ)}
\leanline{4}{(hV\ :\ PotentialAssumptions\ P\ V)\ (n\ :\ ℕ)\ :}
\leanline{4}{(actualProbability\ P\ V).map\ (actualChain\ P\ V\ hV\ n)\ =}
\leanline{6}{(ULA.ulaLawStep\ V\ P.step)\textasciicircum{}[n]\ (ULA.gibbsMeasure\ V)}
\end{leancode}
\leanlabel{How the Lean proof follows the argument}
\texttt{Canonical\allowbreak{}Semigroup.\allowbreak{}state} is the constructed global
    solution. \texttt{state\_\allowbreak{}integral}, \texttt{state\_\allowbreak{}law} and
    \texttt{state\_\allowbreak{}grid} identify its equation, law and grid states.
    \texttt{actual\allowbreak{}Chain\_\allowbreak{}law} identifies the comparison chain with the
    original Gaussian ULA laws.
\end{leanbox}

\section{\texorpdfstring{Proving the actual reversible Markov semigroup}{Proving the actual reversible Markov semigroup}}


Evaluating an observable at \(X_t\) is an isometry from
    \(L^2(\pi)\) into the common probability space. Compress evaluation at
    \(t\) against evaluation at zero to define \(P_t\).
    Future Brownian increments are independent of the complete
    initial/\allowbreak{}Brownian past. Finite shifted Euler maps and their genuine
    \(L^2\) limits prove


    \[\mathbb E[f(X_{a+t})\mid\mathcal F_a]=P_tf(X_a),
    \qquad P_{s+t}=P_sP_t.\]


For reversibility, realize MALA with the same Brownian Gaussian
    proposals and independent uniform accept/\allowbreak{}reject variables. Its anchored
    finite-step laws are exactly reversible. The proved rejection-distance
    estimate and Euler convergence give convergence in probability to the
    actual diffusion; stationarity and density upgrade every observable to
    strong \(L^2\) convergence. Self-adjointness passes to \(P_t\).
    Its contraction and strong continuity follow from actual evaluation
    isometries and state continuity.

\begin{leanbox}{Lean: \texttt{actualEvaluation\_\allowbreak{}markov\_\allowbreak{}pastAt}, \texttt{op\_\allowbreak{}isSelfAdjoint}, \texttt{semigroup}}
\leansource{\texttt{actualEvaluation\_\allowbreak{}markov\_\allowbreak{}pastAt} \quad WassersteinMixing/\allowbreak{}ActualMarkov.lean, lines 120–145}
\begin{leancode}
\leanline{0}{theorem\ actualEvaluation\_markov\_pastAt\ (a\ :\ ℝ≥0)\ (t\ :\ ℝ)\ (ht\ :\ 0\ ≤\ t)}
\leanline{4}{(f\ :\ Lp\ H\ 2\ (ULA.gibbsMeasure\ V))\ :}
\leanline{4}{(condExpL2\ H\ ℝ\ (ULA.lpLangevinPast\_le\ (actualProbability\ P\ V)}
\leanline{6}{(CanonicalSemigroup.initial\ P\ V\ hV)\ (actualBrownianNoise\ P\ V\ hV)\ (a\ :\ ℝ))}
\leanline{6}{(CanonicalSemigroup.evaluation\ ({\leanfallback 𝕜}\ :=\ ℝ)\ P\ V\ hV\ (a\ +\ t)\ f)\ :\ Lp\ H\ 2\ (actualProbability\ P\ V))\ =}
\leanline{6}{CanonicalSemigroup.evaluation\ ({\leanfallback 𝕜}\ :=\ ℝ)\ P\ V\ hV\ a}
\leanline{8}{(CanonicalSemigroup.op\ ({\leanfallback 𝕜}\ :=\ ℝ)\ P\ V\ hV\ t\ f)}
\end{leancode}
\leansource{\texttt{op\_\allowbreak{}isSelfAdjoint} \quad WassersteinMixing/\allowbreak{}SemigroupReversible.lean, lines 70–87}
\begin{leancode}
\leanline{0}{theorem\ op\_isSelfAdjoint\ (P\ :\ Parameters)\ (V\ :\ E\ P.dimension\ →\ ℝ)}
\leanline{4}{(hV\ :\ PotentialAssumptions\ P\ V)\ (t\ :\ ℝ)\ (ht\ :\ 0\ ≤\ t)\ :}
\leanline{4}{IsSelfAdjoint\ (op\ ({\leanfallback 𝕜}\ :=\ {\leanfallback 𝕜})\ (F\ :=\ F)\ P\ V\ hV\ t)}
\end{leancode}
\leansource{\texttt{semigroup} \quad WassersteinMixing/\allowbreak{}SemigroupInstance.lean, lines 55–62}
\begin{leancode}
\leanline{0}{def\ semigroup\ (P\ :\ Parameters)\ (V\ :\ E\ P.dimension\ →\ ℝ)\ (hV\ :\ PotentialAssumptions\ P\ V)\ :}
\leanline{4}{ReversibleSemigroup\ (Lp\ F\ 2\ (ULA.gibbsMeasure\ V))\ where}
\leanline{2}{op\ :=\ op\ P\ V\ hV}
\leanline{2}{zero\ :=\ op\_zero\ P\ V\ hV}
\leanline{2}{add\ :=\ fun\ \{s\ t\}\ hs\ ht\ ↦\ op\_add\ P\ V\ hV\ s\ t\ hs\ ht}
\leanline{2}{selfAdjoint\ :=\ fun\ \{t\}\ ht\ ↦\ op\_isSelfAdjoint\ P\ V\ hV\ t\ ht}
\leanline{2}{norm\_le\ :=\ fun\ \{t\}\ \_\ f\ ↦\ op\_norm\_le\ P\ V\ hV\ t\ f}
\leanline{2}{continuous\ :=\ fun\ f\ ↦\ continuousOn\_op\ P\ V\ hV\ f}
\end{leancode}
\leanlabel{How the Lean proof follows the argument}
\texttt{actual\allowbreak{}Evaluation\_\allowbreak{}markov\_\allowbreak{}past\allowbreak{}At} proves the conditional
    identity. \texttt{op\_\allowbreak{}is\allowbreak{}Self\allowbreak{}Adjoint} uses actual reversible MALA
    approximants. \texttt{Canonical\allowbreak{}Semigroup.\allowbreak{}semigroup} assembles all
    proved fields; no abstract Markov or spectral realization is supplied.
\end{leanbox}

\section{\texorpdfstring{Constructing the bounded Poisson corrector without an eigenbasis}{Constructing the bounded Poisson corrector without an eigenbasis}}


On complex Gibbs observables define
    \(A_hf=\int_0^h(f-P_sf)\,ds\) and
    \(G_N=\sum_{n<N}P_{nh}A_h\).
    Self-adjointness and \(P_t=P_{t/2}^2\) imply
    \(0\preccurlyeq P_t\preccurlyeq P_s\preccurlyeq I\) for
    \(0\le s\le t\). Integration and telescoping give


    \[0\preccurlyeq G_N\preccurlyeq h(I-P_{Nh})\preccurlyeq hI.\]


The increasing bounded operators converge strongly to a genuine
    bounded operator \(G\). Passing the finite telescoping identity to the
    limit gives \((I-P_h)G=A_h\), while
    \(\|G\|\le h\). Commutation and positivity of
    \((h^2I-G^2)(I-P_h)\) give the energy domination


    \[\langle Gf,(I-P_h)Gf\rangle
      \le h^2\langle f,(I-P_h)f\rangle.\]


Pointwise real-to-complex embedding commutes with actual evaluation
    and adjunction. Thus the constructed corrector preserves real
    observables. The extracted real semigroup operator is the actual real
    transition. Set \(g=G\nabla V\).

\begin{leanbox}{Lean: \texttt{corrector\_\allowbreak{}poisson}, \texttt{corrector\_\allowbreak{}energy\_\allowbreak{}le}, \texttt{actualPoissonProcess}}
\leansource{\texttt{corrector\_\allowbreak{}poisson} \quad WassersteinMixing/\allowbreak{}SemigroupPoisson.lean, lines 355–367}
\begin{leancode}
\leanline{0}{theorem\ corrector\_poisson\ (h\ :\ ℝ)\ (hh\ :\ 0\ ≤\ h)\ (x\ :\ H)\ :}
\leanline{4}{S.corrector\ h\ hh\ x\ -\ S.op\ h\ (S.corrector\ h\ hh\ x)\ =\ S.defect\ h\ x}
\end{leancode}
\leansource{\texttt{corrector\_\allowbreak{}energy\_\allowbreak{}le} \quad WassersteinMixing/\allowbreak{}SemigroupPoisson.lean, lines 392–402}
\begin{leancode}
\leanline{0}{theorem\ corrector\_energy\_le\ (h\ :\ ℝ)\ (hh\ :\ 0\ ≤\ h)\ \{s\ :\ ℝ\}\ (hs\ :\ 0\ ≤\ s)\ (x\ :\ H)\ :}
\leanline{4}{inner\ ℝ\ (S.corrector\ h\ hh\ x)}
\leanline{6}{(S.corrector\ h\ hh\ x\ -\ S.op\ s\ (S.corrector\ h\ hh\ x))\ ≤}
\leanline{6}{h\ \textasciicircum{}\ 2\ *\ inner\ ℝ\ x\ (x\ -\ S.op\ s\ x)}
\end{leancode}
\leansource{\texttt{actualPoissonProcess} \quad WassersteinMixing/\allowbreak{}ActualPoisson.lean, lines 81–88}
\begin{leancode}
\leanline{0}{def\ actualPoissonProcess\ (P\ :\ Parameters)\ (V\ :\ E\ P.dimension\ →\ ℝ)}
\leanline{4}{(hV\ :\ PotentialAssumptions\ P\ V)\ :}
\leanline{4}{PoissonProcess\ (Lp\ (E\ P.dimension)\ 2\ (ULA.gibbsMeasure\ V))\ (E\ P.dimension)}
\leanline{6}{(actualProbability\ P\ V)\ P.step\ P.step\_pos\ :=}
\leanline{2}{semigroupPoissonProcess\ (actualProbability\ P\ V)\ (actualObservableSemigroup\ P\ V\ hV)}
\leanline{4}{(actualRealEmbedding\ P\ V)\ (actualObservableSemigroup\_preservesReal\ P\ V\ hV)}
\leanline{4}{P.step\ P.step\_pos\ (actualFiltration\ P\ V\ hV)\ (actualGridEvaluation\ P\ V\ hV)}
\leanline{4}{(actualGridEvaluation\_adapted\ P\ V\ hV)\ (actualGridEvaluation\_markov\ P\ V\ hV)}
\end{leancode}
\leanlabel{How the Lean proof follows the argument}
\texttt{corrector\_\allowbreak{}norm\_\allowbreak{}le},
    \texttt{corrector\_\allowbreak{}poisson} and \texttt{corrector\_\allowbreak{}energy\_\allowbreak{}le}
    prove the required conclusions from the constructed semigroup.
    \texttt{actual\allowbreak{}Poisson\allowbreak{}Process} transfers them to the actual adapted
    grid evaluations.
\end{leanbox}

\section{\texorpdfstring{Deriving the actual score increment and defect bounds}{Deriving the actual score increment and defect bounds}}


Gibbs integration gives
    \(\|\nabla V\|_{L^2(\pi)}^2\le\beta d\).
    The actual integral equation and Brownian second moment imply


    \[\|X_{kh+s}-X_{kh}\|_{L^2}^2
      \le d(2s^2\beta+4s).\]


Gradient Lipschitzness and \(s\beta\le1\) then give


    \[\|\nabla V(X_{kh+s})-\nabla V(X_{kh})\|_{L^2}^2
      \le6s\beta^2d.\]


The true local defect is
    \(D_k=\int_0^h(\nabla V(X_{kh})-\nabla V(X_{kh+s}))\,ds\).
    Integral Cauchy–Schwarz gives
    \(\|D_k\|_{L^2}^2\le3h^3\beta^2d\).
    The corrector bound yields
    \(\|g(X_{kh})\|_{L^2}^2\le h^2\beta d\), and its energy bound
    gives a corrector-increment budget \(6h^3\beta^2d\).

\begin{leanbox}{Lean: \texttt{actualObservableScore\_\allowbreak{}energy}, \texttt{actualScoreIncrement\_\allowbreak{}energy}, \texttt{actualDefect\_\allowbreak{}energy}}
\leansource{\texttt{actualObservableScore\_\allowbreak{}energy} \quad WassersteinMixing/\allowbreak{}ActualScoreObservable.lean, lines 64–69}
\begin{leancode}
\leanline{0}{theorem\ actualObservableScore\_energy\ :}
\leanline{4}{‖actualObservableScore\ P\ V\ hV‖\ \textasciicircum{}\ 2\ ≤\ P.beta\ *\ P.dimension}
\end{leancode}
\leansource{\texttt{actualScoreIncrement\_\allowbreak{}energy} \quad WassersteinMixing/\allowbreak{}ActualLocalError.lean, lines 117–126}
\begin{leancode}
\leanline{0}{theorem\ actualScoreIncrement\_energy\ (n\ :\ ℕ)\ (s\ :\ ℝ)\ (hs\ :\ s\ ∈\ Icc\ 0\ P.step)\ :}
\leanline{4}{‖actualScoreIncrement\ P\ V\ hV\ n\ s‖\ \textasciicircum{}\ 2\ ≤\ 6\ *\ s\ *\ P.beta\ \textasciicircum{}\ 2\ *\ P.dimension}
\end{leancode}
\leansource{\texttt{actualDefect\_\allowbreak{}energy} \quad WassersteinMixing/\allowbreak{}ActualLocalError.lean, lines 154–158}
\begin{leancode}
\leanline{0}{theorem\ actualDefect\_energy\ (n\ :\ ℕ)\ :}
\leanline{4}{‖actualDefect\ P\ V\ hV\ n‖\ \textasciicircum{}\ 2\ ≤\ 3\ *\ P.step\ \textasciicircum{}\ 3\ *\ P.beta\ \textasciicircum{}\ 2\ *\ P.dimension}
\end{leancode}
\leanlabel{How the Lean proof follows the argument}
\texttt{actual\allowbreak{}Observable\allowbreak{}Score\_\allowbreak{}energy},
    \texttt{actual\allowbreak{}Score\allowbreak{}Increment\_\allowbreak{}energy} and
    \texttt{actual\allowbreak{}Defect\_\allowbreak{}energy} derive these constants for the genuine
    diffusion. They are not inputs to the final theorem.
\end{leanbox}

\section{\texorpdfstring{Identifying and centering the actual martingale residual}{Identifying and centering the actual martingale residual}}


The proved conditional Markov identity and commutation with
    the Bochner time integral identify the actual defect mean:


    \[\mathbb E[D_k\mid\mathcal F_{kh}]
      =\int_0^h(I-P_s)\nabla V(X_{kh})\,ds
      =(I-P_h)g(X_{kh}).\]


Hence
    \(M_{k+1}=D_k-[g(X_{kh})-g(X_{(k+1)h})]\) is conditionally
    centered. The separate budgets \(3h^3\beta^2d\) and
    \(6h^3\beta^2d\), together with
    \(\|u-v\|^2\le2\|u\|^2+2\|v\|^2\), give


    \[\|M_{k+1}\|_{L^2}^2\le18h^3\beta^2d.\]


The predictable shifted drift belongs to the actual Brownian past,
    so conditional orthogonality removes its cross term with \(M_{k+1}\).

\begin{leanbox}{Lean: \texttt{actualDefect\_\allowbreak{}conditionalMean}, \texttt{martingaleAt\_\allowbreak{}energy\_\allowbreak{}le}, \texttt{shifted\_\allowbreak{}energy\_\allowbreak{}step}}
\leansource{\texttt{actualDefect\_\allowbreak{}conditionalMean} \quad WassersteinMixing/\allowbreak{}ActualDefectMean.lean, lines 81–92}
\begin{leancode}
\leanline{0}{theorem\ actualDefect\_conditionalMean\ (n\ :\ ℕ)\ :}
\leanline{4}{(condExpL2\ (E\ P.dimension)\ ℝ\ (actualPast\_le\ P\ V\ hV\ n)}
\leanline{6}{(actualDefect\ P\ V\ hV\ n)\ :\ Lp\ (E\ P.dimension)\ 2\ (actualProbability\ P\ V))\ =}
\leanline{6}{actualGridEvaluation\ P\ V\ hV\ n}
\leanline{8}{(∫\ s\ in\ 0..P.step,\ actualObservableScore\ P\ V\ hV\ -}
\leanline{10}{CanonicalSemigroup.op\ ({\leanfallback 𝕜}\ :=\ ℝ)\ P\ V\ hV\ s\ (actualObservableScore\ P\ V\ hV))}
\end{leancode}
\leansource{\texttt{martingaleAt\_\allowbreak{}energy\_\allowbreak{}le} \quad WassersteinMixing/\allowbreak{}PoissonProcess.lean, lines 148–156}
\begin{leancode}
\leanline{0}{theorem\ martingaleAt\_energy\_le\ (f\ :\ H)\ (D\ :\ ℕ\ →\ Lp\ E\ 2\ Q)\ (β\ d\ :\ ℝ)\ (n\ :\ ℕ)}
\leanline{4}{(hdefect\ :\ ‖D\ n‖\ \textasciicircum{}\ 2\ ≤\ 3\ *\ h\ \textasciicircum{}\ 3\ *\ β\ \textasciicircum{}\ 2\ *\ d)}
\leanline{4}{(hincrement\ :\ ‖S.eval\ (n\ +\ 1)\ f\ -\ S.eval\ n\ f‖\ \textasciicircum{}\ 2\ ≤\ 6\ *\ h\ *\ β\ \textasciicircum{}\ 2\ *\ d)\ :}
\leanline{4}{‖S.martingaleAt\ f\ D\ n‖\ \textasciicircum{}\ 2\ ≤\ 18\ *\ h\ \textasciicircum{}\ 3\ *\ β\ \textasciicircum{}\ 2\ *\ d}
\end{leancode}
\leansource{\texttt{shifted\_\allowbreak{}energy\_\allowbreak{}step} \quad WassersteinMixing/\allowbreak{}PoissonComparison.lean, lines 134–148}
\begin{leancode}
\leanline{0}{theorem\ shifted\_energy\_step\ (n\ :\ ℕ)\ :}
\leanline{4}{‖C.error\ (n\ +\ 1)\ +\ S.correctorAt\ score\ (n\ +\ 1)‖\ \textasciicircum{}\ 2\ ≤}
\leanline{6}{P.contraction\ *\ ‖C.error\ n\ +\ S.correctorAt\ score\ n‖\ \textasciicircum{}\ 2\ +\ P.energyForcing}
\end{leancode}
\leanlabel{How the Lean proof follows the argument}
\texttt{actual\allowbreak{}Defect\_\allowbreak{}conditional\allowbreak{}Mean} identifies the true
    integral defect. The Poisson process derives martingale centering and
    its energy bound; the synchronous comparison then derives the shifted
    energy recurrence.
\end{leanbox}

\section{\texorpdfstring{Using the genuine averaged Hessian and shifted energy}{Using the genuine averaged Hessian and shifted energy}}


The literal segment average of the Hessian between the actual
    diffusion and ULA states is symmetric, lies between \(\alpha I\) and
    \(\beta I\), and maps their difference to the true gradient difference.
    Writing \(Z_k=X_{kh}-\widehat X_k\) and \(U_k=Z_k+g(X_{kh})\),
    the exact synchronous updates give


    \[U_{k+1}=(I-hA_k)U_k+hA_kg(X_{kh})+M_{k+1}.\]


The positive contraction inequality yields


    \[\|(I-hA_k)u+hA_kv\|^2
      \le(1-\alpha h)\|u\|^2+h\beta\|v\|^2.\]


After conditional orthogonality, the forcing constant is
    \(1+18=19\):


    \[\|U_{k+1}\|_{L^2}^2
      \le(1-\alpha h)\|U_k\|_{L^2}^2+19h^3\beta^2d.\]


Starting the two processes together gives the uniform floor
    \(19h^2\kappa\beta d\). The triangle inequality, with
    \(\sqrt{19}\le5\), gives
    \(\|Z_k\|_{L^2}\le6h\sqrt{\kappa\beta d}\).

\begin{leanbox}{Lean: \texttt{actual\_\allowbreak{}curvature\_\allowbreak{}error\_\allowbreak{}step}, \texttt{curvature\_\allowbreak{}weighted\_\allowbreak{}interpolation}, \texttt{shifted\_\allowbreak{}norm\_\allowbreak{}le\_\allowbreak{}bias}}
\leansource{\texttt{actual\_\allowbreak{}curvature\_\allowbreak{}error\_\allowbreak{}step} \quad WassersteinMixing/\allowbreak{}ActualCurvature.lean, lines 25–37}
\begin{leancode}
\leanline{0}{theorem\ actual\_curvature\_error\_step\ (n\ :\ ℕ)\ :}
\leanline{4}{actualReference\ P\ V\ hV\ (n\ +\ 1)\ -\ actualChain\ P\ V\ hV\ (n\ +\ 1)\ =}
\leanline{6}{(actualReference\ P\ V\ hV\ n\ -\ actualChain\ P\ V\ hV\ n)\ -}
\leanline{8}{P.step\ •\ actualCurvature\ P\ V\ hV\ n\ (actualReference\ P\ V\ hV\ n\ -\ actualChain\ P\ V\ hV\ n)\ +}
\leanline{10}{actualDefect\ P\ V\ hV\ n}
\end{leancode}
\leansource{\texttt{curvature\_\allowbreak{}weighted\_\allowbreak{}interpolation} \quad WassersteinMixing/\allowbreak{}ShiftedEnergy.lean, lines 72–93}
\begin{leancode}
\leanline{0}{theorem\ curvature\_weighted\_interpolation\ (P\ :\ Parameters)\ (A\ :\ H\ →L[ℝ]\ H)}
\leanline{4}{(hsym\ :\ ∀\ x\ y,\ inner\ ℝ\ (A\ x)\ y\ =\ inner\ ℝ\ x\ (A\ y))}
\leanline{4}{(hlower\ :\ ∀\ x,\ P.alpha\ *\ ‖x‖\ \textasciicircum{}\ 2\ ≤\ inner\ ℝ\ x\ (A\ x))}
\leanline{4}{(hupper\ :\ ∀\ x,\ inner\ ℝ\ x\ (A\ x)\ ≤\ P.beta\ *\ ‖x‖\ \textasciicircum{}\ 2)\ (u\ v\ :\ H)\ :}
\leanline{4}{‖u\ -\ P.step\ •\ A\ u\ +\ P.step\ •\ A\ v‖\ \textasciicircum{}\ 2\ ≤}
\leanline{6}{P.contraction\ *\ ‖u‖\ \textasciicircum{}\ 2\ +\ P.step\ *\ P.beta\ *\ ‖v‖\ \textasciicircum{}\ 2}
\end{leancode}
\leansource{\texttt{shifted\_\allowbreak{}norm\_\allowbreak{}le\_\allowbreak{}bias} \quad WassersteinMixing/\allowbreak{}Recurrence.lean, lines 88–121}
\begin{leancode}
\leanline{0}{theorem\ shifted\_norm\_le\_bias\ \{H\ :\ Type*\}\ [NormedAddCommGroup\ H]}
\leanline{4}{(P\ :\ Parameters)\ (z\ g\ :\ H)}
\leanline{4}{(hu\ :\ ‖z\ +\ g‖\ \textasciicircum{}\ 2\ ≤\ P.energyFloor)}
\leanline{4}{(hg\ :\ ‖g‖\ \textasciicircum{}\ 2\ ≤\ P.correctorBudget)\ :\ ‖z‖\ ≤\ P.bias}
\end{leancode}
\leanlabel{How the Lean proof follows the argument}
\texttt{actual\_\allowbreak{}curvature\_\allowbreak{}error\_\allowbreak{}step} uses the genuine
    averaged-Hessian secant identity.
    \texttt{curvature\_\allowbreak{}weighted\_\allowbreak{}interpolation} proves the deterministic
    inequality. \texttt{shifted\_\allowbreak{}energy\_\allowbreak{}uniform} and
    \texttt{shifted\_\allowbreak{}norm\_\allowbreak{}le\_\allowbreak{}bias} derive the floor and literal
    coefficient six.
\end{leanbox}

\section{\texorpdfstring{Constructing the actual Gaussian ULA invariant law}{Constructing the actual Gaussian ULA invariant law}}


Let \(F(x)=x-h\nabla V(x)\), whose Lipschitz factor is
    \(q=1-\alpha h<1\). On an iid Gaussian past
    \(\omega=(\omega_0,\omega_1,\ldots)\), let \(\tau\) drop its first entry.
    The tail preserves the product law and is independent of \(\omega_0\).
    On the complete \(L^2\) space define


    \[\Phi(X)(\omega)=F(X(\tau\omega))+\sqrt{2h}\,\omega_0.\]


This is a contraction with factor \(q\), so Banach's theorem
    constructs a square-integrable fixed point. Its law \(\widehat\pi\)
    is a probability measure with finite second moment. The actual
    recursion and head/\allowbreak{}tail independence prove invariance under the
    original independent-Gaussian ULA transition. The proved two-law
    contraction gives
    \(W_2(T^k\mu,\widehat\pi)\le q^kW_2(\mu,\widehat\pi)\to0\).

\begin{leanbox}{Lean: \texttt{invariantLaw\_\allowbreak{}invariant}, \texttt{invariantLaw\_\allowbreak{}attraction}}
\leansource{\texttt{invariantLaw\_\allowbreak{}invariant} \quad WassersteinMixing/\allowbreak{}Invariant.lean, lines 236–242}
\begin{leancode}
\leanline{0}{theorem\ invariantLaw\_invariant\ (P\ :\ Parameters)\ (V\ :\ E\ →\ ℝ)}
\leanline{4}{(hV\ :\ PotentialAssumptions\ P\ V)\ :}
\leanline{4}{ULA.ulaLawStep\ V\ P.step\ (invariantLaw\ P\ V\ hV)\ =\ invariantLaw\ P\ V\ hV}
\end{leancode}
\leansource{\texttt{invariantLaw\_\allowbreak{}attraction} \quad WassersteinMixing/\allowbreak{}Invariant.lean, lines 256–265}
\begin{leancode}
\leanline{0}{theorem\ invariantLaw\_attraction\ (P\ :\ Parameters)\ (V\ :\ E\ →\ ℝ)}
\leanline{4}{(hV\ :\ PotentialAssumptions\ P\ V)\ (μ\ :\ Measure\ E)\ [IsProbabilityMeasure\ μ]}
\leanline{4}{(hμ\ :\ MemLp\ (id\ :\ E\ →\ E)\ 2\ μ)\ :}
\leanline{4}{Filter.Tendsto\ (fun\ k\ :\ ℕ\ ↦}
\leanline{6}{ULA.wasserstein2Real\ ((ULA.ulaLawStep\ V\ P.step)\textasciicircum{}[k]\ μ)\ (invariantLaw\ P\ V\ hV))}
\leanline{6}{Filter.atTop\ ({\leanfallback 𝓝}\ 0)}
\end{leancode}
\leanlabel{How the Lean proof follows the argument}
\texttt{invariant\allowbreak{}Law} specializes the genuine random
    fixed-point construction. \texttt{invariant\allowbreak{}Law\_\allowbreak{}invariant} and
    \texttt{invariant\allowbreak{}Law\_\allowbreak{}attraction} prove the exact stationary and
    limiting properties needed in Theorem 1.
\end{leanbox}

\section{\texorpdfstring{Passing the actual stationary comparison to transport}{Passing the actual stationary comparison to transport}}


The synchronous trajectories give actual couplings of
    \(\pi\) and each stationary-start Gaussian ULA law. Their uniform
    \(L^2\) bound therefore bounds the actual quadratic transport
    infimum. To prove its triangle inequality, disintegrate two couplings
    with the same middle marginal and construct a genuine triple law.
    Minkowski bounds its endpoint cost; taking both coupling infima gives


    \[W_2(\mu,\xi)\le W_2(\mu,\nu)+W_2(\nu,\xi).\]


Apply this with the stationary-start ULA law between \(\pi\) and
    \(\widehat\pi\), then use the proved attraction to the constructed
    invariant law. Theorem 1 follows. Combining it with actual ULA
    contraction gives the finite-time inequality of Corollary 1.1.
    Second moments prove finiteness before distances are converted to
    real numbers.

\begin{leanbox}{Lean: \texttt{stationary\_\allowbreak{}transport\_\allowbreak{}bound}, \texttt{wasserstein2Real\_\allowbreak{}triangle}, \texttt{ula\_\allowbreak{}iterate\_\allowbreak{}wasserstein2Real\_\allowbreak{}le}}
\leansource{\texttt{stationary\_\allowbreak{}transport\_\allowbreak{}bound} \quad WassersteinMixing/\allowbreak{}PoissonComparison.lean, lines 163–167}
\begin{leancode}
\leanline{0}{theorem\ stationary\_transport\_bound\ [BorelSpace\ E]\ [SecondCountableTopology\ E]\ (n\ :\ ℕ)\ :}
\leanline{4}{ULA.wasserstein2Real\ (Q.map\ (C.reference\ 0))\ (Q.map\ (C.chain\ n))\ ≤\ P.bias}
\end{leancode}
\leansource{\texttt{wasserstein2Real\_\allowbreak{}triangle} \quad WassersteinMixing/\allowbreak{}TransportGeometry.lean, lines 175–181}
\begin{leancode}
\leanline{0}{theorem\ wasserstein2Real\_triangle\ (μ\ ν\ ξ\ :\ Measure\ H)}
\leanline{4}{[IsProbabilityMeasure\ μ]\ [IsProbabilityMeasure\ ν]\ [IsProbabilityMeasure\ ξ]}
\leanline{4}{(hμν\ :\ ULA.wasserstein2\ μ\ ν\ ≠\ ⊤)\ (hνξ\ :\ ULA.wasserstein2\ ν\ ξ\ ≠\ ⊤)\ :}
\leanline{4}{ULA.wasserstein2Real\ μ\ ξ\ ≤\ ULA.wasserstein2Real\ μ\ ν\ +\ ULA.wasserstein2Real\ ν\ ξ}
\end{leancode}
\leansource{\texttt{ula\_\allowbreak{}iterate\_\allowbreak{}wasserstein2Real\_\allowbreak{}le} \quad WassersteinMixing/\allowbreak{}Transport.lean, lines 109–121}
\begin{leancode}
\leanline{0}{theorem\ ula\_iterate\_wasserstein2Real\_le\ (P\ :\ Parameters)\ (V\ :\ H\ →\ ℝ)}
\leanline{4}{(hV\ :\ PotentialAssumptions\ P\ V)\ (μ\ ν\ :\ Measure\ H)}
\leanline{4}{[IsProbabilityMeasure\ μ]\ [IsProbabilityMeasure\ ν]}
\leanline{4}{(hfinite\ :\ ULA.wasserstein2\ μ\ ν\ ≠\ ⊤)\ (k\ :\ ℕ)\ :}
\leanline{4}{ULA.wasserstein2Real\ ((ULA.ulaLawStep\ V\ P.step)\textasciicircum{}[k]\ μ)}
\leanline{6}{((ULA.ulaLawStep\ V\ P.step)\textasciicircum{}[k]\ ν)\ ≤}
\leanline{6}{P.contraction\ \textasciicircum{}\ k\ *\ ULA.wasserstein2Real\ μ\ ν}
\end{leancode}
\leanlabel{How the Lean proof follows the argument}
\texttt{stationary\_\allowbreak{}transport\_\allowbreak{}bound} uses the actual
    synchronous coupling. \texttt{glue\allowbreak{}Couplings} constructs the
    triple law and \texttt{wasserstein2\allowbreak{}Real\_\allowbreak{}triangle} proves the
    finite-distance triangle inequality. The final theorem instantiates
    the comparison from the source Hessian assumptions.
\end{leanbox}

\section{\texorpdfstring{Recovering the exact scale-invariant iteration count}{Recovering the exact scale-invariant iteration count}}


For \(d>0\) and \(\varepsilon>0\), write


    \[m=\max\left(1,\frac{12\sqrt d}{\varepsilon}\right),\qquad
    h=\frac1{\beta m},\qquad K=\frac\beta\alpha m.\]


This is the source minimum formula and gives \(\alpha hK=1\).
    Since
    \(\sqrt\alpha\sqrt{(\beta/\alpha)\beta d}=\beta\sqrt d\),
    the scaled bias is at most \(\varepsilon/2\).
    For \(t=\alpha h\in(0,1]\), use \((1-t)^k\le e^{-tk}\).
    The displayed source iteration restriction makes the scaled transient
    at most \(\varepsilon/2\). Adding the two budgets proves the
    literal accuracy conclusion, including zero initial error.

\begin{leanbox}{Lean: \texttt{paperStep\_\allowbreak{}eq\_\allowbreak{}inv}, \texttt{paper\_\allowbreak{}complexity\_\allowbreak{}bound}, \texttt{accuracy\_\allowbreak{}bound}}
\leansource{\texttt{paperStep\_\allowbreak{}eq\_\allowbreak{}inv} \quad WassersteinMixing/\allowbreak{}Complexity.lean, lines 27–40}
\begin{leancode}
\leanline{0}{theorem\ paperStep\_eq\_inv\ (β\ d\ ε\ :\ ℝ)\ (hβ\ :\ 0\ <\ β)\ (hd\ :\ 0\ <\ d)\ (hε\ :\ 0\ <\ ε)\ :}
\leanline{4}{paperStep\ β\ d\ ε\ =\ 1\ /\ (β\ *\ max\ 1\ (12\ *\ Real.sqrt\ d\ /\ ε))}
\end{leancode}
\leansource{\texttt{paper\_\allowbreak{}complexity\_\allowbreak{}bound} \quad WassersteinMixing/\allowbreak{}Complexity.lean, lines 118–142}
\begin{leancode}
\leanline{0}{theorem\ paper\_complexity\_bound\ (α\ β\ d\ R\ ε\ :\ ℝ)\ (hα\ :\ 0\ <\ α)\ (hαβ\ :\ α\ ≤\ β)}
\leanline{4}{(hd\ :\ 0\ <\ d)\ (hR\ :\ 0\ ≤\ R)\ (hε\ :\ 0\ <\ ε)\ (k\ :\ ℕ)}
\leanline{4}{(hk\ :\ iterationPrefactor\ α\ β\ d\ ε\ *}
\leanline{6}{Real.log\ (1\ +\ 2\ *\ Real.sqrt\ α\ *\ R\ /\ ε)\ ≤\ (k\ :\ ℝ))\ :}
\leanline{4}{Real.sqrt\ α\ *\ ((1\ -\ α\ *\ paperStep\ β\ d\ ε)\ \textasciicircum{}\ k\ *\ R\ +}
\leanline{6}{6\ *\ paperStep\ β\ d\ ε\ *\ Real.sqrt\ ((β\ /\ α)\ *\ β\ *\ d))\ ≤\ ε}
\end{leancode}
\leansource{\texttt{accuracy\_\allowbreak{}bound} \quad WassersteinMixing/\allowbreak{}Theorem.lean, lines 44–54}
\begin{leancode}
\leanline{0}{theorem\ accuracy\_bound\ (P\ :\ Parameters)\ (V\ :\ E\ P.dimension\ →\ ℝ)}
\leanline{4}{(hA\ :\ HessianAssumptions\ P\ V)\ (μ\ :\ Measure\ (E\ P.dimension))\ [IsProbabilityMeasure\ μ]}
\leanline{4}{(hμ\ :\ MemLp\ (id\ :\ E\ P.dimension\ →\ E\ P.dimension)\ 2\ μ)\ (ε\ :\ ℝ)\ (hε\ :\ 0\ <\ ε)}
\leanline{4}{(hstep\ :\ P.step\ =\ paperStep\ P.beta\ P.dimension\ ε)\ (n\ :\ ℕ)}
\leanline{4}{(hn\ :\ iterationPrefactor\ P.alpha\ P.beta\ P.dimension\ ε\ *}
\leanline{6}{Real.log\ (1\ +\ 2\ *\ Real.sqrt\ P.alpha\ *}
\leanline{8}{ULA.wasserstein2Real\ μ\ (invariantLaw\ P\ V\ hA.toPotentialAssumptions)\ /\ ε)\ ≤\ (n\ :\ ℝ))\ :}
\leanline{4}{Real.sqrt\ P.alpha\ *}
\leanline{6}{ULA.wasserstein2Real\ ((ULA.ulaLawStep\ V\ P.step)\textasciicircum{}[n]\ μ)\ (ULA.gibbsMeasure\ V)\ ≤\ ε}
\end{leancode}
\leanlabel{How the Lean proof follows the argument}
\texttt{paper\allowbreak{}Step\_\allowbreak{}eq\_\allowbreak{}inv} identifies the two step
    formulas. \texttt{paper\_\allowbreak{}complexity\_\allowbreak{}bound} proves their exact
    scalar consequence. \texttt{accuracy\_\allowbreak{}bound} applies it to the
    actual Gaussian chain and constructed invariant law.
\end{leanbox}

\end{document}
